\documentclass[
paper=a4,
12pt,
twoside=off,
bibliography=totoc,
listof=totoc,
appendixprefix,
headings=small]{scrbook} 

\usepackage[
paper=a4paper,
left=15mm,
right=25mm,
top=25mm,
bottom=25mm,
bindingoffset=10mm]{geometry}
\usepackage{scrhack}

\usepackage[utf8]{inputenc}
\usepackage[greek,english]{babel}

\usepackage{setspace}
\onehalfspacing

\usepackage{lmodern}
\renewcommand*\familydefault{\sfdefault} % base font is to be sans serif
\usepackage[T1]{fontenc}
\usepackage{microtype} % für die Mikrotypografie (besseres Schriftbild)

\usepackage{url} % URL's formatieren (z.B. in Literatur) 
\usepackage[hidelinks]{hyperref} % für Hyperlinks in PDF-Dokumenten

\usepackage{longtable}
\usepackage{multicol, multirow}

\usepackage{cellspace}
\setlength{\tabcolsep}{12pt} %gap between text and border
\renewcommand{\arraystretch}{1.2} %height of cells
\setlength{\arrayrulewidth}{1pt}



\begin{document}

bbbbbbbb


\begin{longtable}[]{|m{1.5cm}|m{3.6cm}m{5cm}m{2.2cm}|}
    \caption{Composition of Buffers used in this thesis}
    \label{tab:buffers_used}
    \\

    \hline
        \textbf{Name} & \multicolumn{3}{c|}{\textbf{Components}} \\
    \hline
    \endfirsthead

    \multirow{4}{4em}{\textbf{Elution Buffer}} 
        & 20 mM & N-2-hydroxyethylpiperazine-N'-2-ethane sulphonic acid & by neoFroxx \\
        & 50/100/250/500 mM & imidazole & by Thermo Scientific \\
        & 300 mM & sodium chloride & by neoFroxx \\
        & \multicolumn{3}{c|}{pH adjusted to 8.0} \\    
    \hline
    
    \multirow{4}{4em}{\textbf{Elution Buffer}} 
        & 20 mM & 4-(2-Hydroxyethyl)piperazine-1-ethanesulfonic acid & by neoFroxx \\
        & 50/100/250/500 mM & imidazole & by Thermo Scientific \\
        & 300 mM & sodium chloride & by neoFroxx \\
        & \multicolumn{3}{c|}{pH adjusted to 8.0} \\    
    \hline
    
\end{longtable}


\end{document}